\section{Introduction}\label{sec:intro}

\subsection{The problem}
Many files must be kept but are rarely opened: finished projects, scans, archives, old
photographs. Keeping them on the computer costs space; moving them to a storage service hands
them to a provider that can read them, or that holds the keys that would let it. Many people who
write software already have something else: private GitHub repositories, reached through a
well-documented HTTPS interface, with an access-controlled place to put data.

Two facts stand in the way of simply putting files there. A private repository is private only
by GitHub's access control: GitHub, every collaborator, and anyone holding a token with access
can read every byte of it, and so can everyone if the repository is ever made public. And GitHub
is explicit that ``Git is not designed to serve as a backup tool''~\cite{gh-large}; a repository
can be deleted, an account suspended, a token revoked.

\subsection{The idea}
GhostVault treats the repositories as an untrusted store of opaque objects. A file is
\emph{ghosted}: it is encrypted on the computer, uploaded in chunks, checked on GitHub, and then
removed from the computer. What remains locally is a \emph{ghost}: a small record that holds the
file's key and the location of every chunk, shown as a card in GhostVault's window.
\emph{Restoring} a ghost downloads and checks its chunks, writes the file into a folder the user
chooses, and then removes its copy from GitHub. Everything GitHub receives about a file is
encrypted: it never sees the file's key, its name or its content.

The consequence that shapes the whole design is that GhostVault deletes the user's files on
purpose. Encryption is the easier part; the harder part is deciding, after crashes, network
failures, misleading replies and other programs touching the same files, when a deletion is
allowed. The original is deleted only once its encrypted copy has been shown to be on GitHub, and
the remote copy only once the restored file is safely on disk. Most of this report is about how
those two decisions are made and recorded.

\subsection{What this report covers}
This is a systems report about GhostVault 1.0.0. It describes what the system is made of, the
formats it writes, the rules its operations follow and why, and how it was checked. Every
statement about behaviour was checked against the source code; every number comes from a
measurement whose method is given, run on one stated machine; and what has not been verified is
listed in \secref{sec:status}. The report aims to be as precise about what GhostVault does not do
as about what it does.

\secref{sec:overview} gives an overview and \secref{sec:principles} the principles behind the
design. \secref{sec:threats} sets out the threat model and the boundaries of the protection.
Sections~\ref{sec:crypto} to \ref{sec:recovery} describe the cryptography, the local state and
its journal, ghosting, restoring, the use of GitHub, and recovery.
Sections~\ref{sec:app} and~\ref{sec:impl} cover the desktop application, its standalone builds and
the implementation. Sections~\ref{sec:verification} and~\ref{sec:performance} report verification
and performance, and \secref{sec:limitations} the limitations.
